% Any commands that need to be run before the \documentclass command below
% should be done in the preclass file below.
% Occasionally, some commands need to be run before the documentclass
% declaration. Place them here.

% Copied from the ACM sigconf template.
%% Commands for TeXCount
%TC:macro \cite [option:text,text]
%TC:macro \citep [option:text,text]
%TC:macro \citet [option:text,text]
%TC:envir table 0 1
%TC:envir table* 0 1
%TC:envir tabular [ignore] word
%TC:envir displaymath 0 word
%TC:envir math 0 word
%TC:envir comment 0 0
%%
%% The tagged file should start with the metadata commands.
\DocumentMetadata{
 pdfversion=2.0,pdfstandard=ua-2,
 testphase=latest,
 tagging=on,
 tagging-setup={math/setup=mathml-SE},
 lang = en
 }


% Read your conference/journal submission guidelines for what options to include
% here.
\documentclass[sigconf]{acmart}

% Do any other setup of packages, commands, in the relevant setup file, imported
% with the following file.
% This is the setup file. It contains usepackage commands, acronyms, and other
% useful commands needed to set up the paper. Uncomment the relevant sections to
% use them.

% Some commands rely on a conditional to have different behavior at submit time
% (e.g., any author comments are stripped from submission versions of the final
% paper). Uncomment \submittrue below to use submit-time behavior, and
% \submitfalse to use default behavior. If both lines are uncommented, whatever
% comes later will take precedence.
\newif\ifsubmit{}  % Flag for distinguishing submission commands
%\submittrue{}
\submitfalse{}

% Similarly, the cameraready conditional below provides for different behavior
% in a camera-ready version, like real author affiliations.
\newif\ifcameraready{}  % Flag for distinguishing submission commands
%\camerareadytrue{}
\camerareadyfalse{}

% Setup preamble
%
% This is the place for anything that needs to be executed before package
% includes, command definitions, and the like. If copying from a template,
% anything in the template that doesn't fit elsewhere in setup and comes before
% the \begin{document} command should go here.

% For example, in ACM paperes, there's often a copyright section like the one
% below. Uncomment it 
%\setcopyright{acmlicensed}
%\copyrightyear{2018}
%\acmYear{2018}
%\acmDOI{XXXXXXX.XXXXXXX}
%\acmConference[Conference acronym 'XX]{Make sure to enter the correct
  %conference title from your rights confirmation email}{June 03--05,
  %2018}{Woodstock, NY}
%\acmISBN{978-1-4503-XXXX-X/2018/06}


% Declare any packages and options you are using here.

% In the ACM template, there is a list of allowed packages and list of packages
% already included as part of acmart. See here:
% https://authors.acm.org/proceedings/production-information/accepted-latex-packages

% INCLUDED PACKAGES %%%%%%%%%%%%%%%%%%%%%%%%%%%%%%%%%%%%%%%%%%%%%%%%%%%%%%%%%%%%
% The acmart template includes some packages already. DO NOT UNCOMMENT ANYTHING
% IN THESE LISTS - they're just there to help you remember which ones you're
% using that are included with acmart.
% Enable features useful for typesetting mathematical notation (e.g., multiline
% alignment for equations).
%\usepackage{amsmath}
% Balance the final page columns.
%\usepackage{balance}
% Make tables look nicer.
%\usepackage{booktabs}
% Allow figures
%\usepackage{graphicx}
% Create hypertext links for citations and references.
%\usepackage{hyperref}
% Define colors for comments, etc.
%\usepackage{xcolor}
%%%%%%%%%%%%%%%%%%%%%%%%%%%%%%%%%%%%%%%%%%%%%%%%%%%%%%%%%%%%%%%%%%%%%%%%%%%%%%%%

% ACTIVE PACKAGES %%%%%%%%%%%%%%%%%%%%%%%%%%%%%%%%%%%%%%%%%%%%%%%%%%%%%%%%%%%%%%
% Used packages. Keep packages in alphabetical order, except for hyperref below.
%
% Use acronyms (define actual acronyms in setup/acro.tex). The printonlyused
% option is probably not useful in papers since an acronym list isn't printed at
% all, but leave it here just in case.
\usepackage[printonlyused]{acronym}
% Highly customizable enumeration, including inline lists.
\usepackage[inline]{enumitem}
% For spacing after user-defined commands and others.
\usepackage{xspace}
%%%%%%%%%%%%%%%%%%%%%%%%%%%%%%%%%%%%%%%%%%%%%%%%%%%%%%%%%%%%%%%%%%%%%%%%%%%%%%%%

% OTHER PACKAGES %%%%%%%%%%%%%%%%%%%%%%%%%%%%%%%%%%%%%%%%%%%%%%%%%%%%%%%%%%%%%%%
% Packages that are not used, but may be useful in certain situations. Move and
% uncomment to the list above to use.
%
% Use acronyms (define actual acronyms in setup/acro.tex).
%\usepackage{acronym}
%
% Enable algorithm descriptions in pseudocode.
%\usepackage{algorithmic}
%
% Enable features useful for typesetting mathematical notation (e.g., multiline
% alignment for equations).
%\usepackage{amsmath}
%
% Load AMS fonts.
%\usepackage{amsfonts}
%
% Enable commands for certain mathematical symbols (e.g., \varnothing for the
% empty set).
%\usepackage{amssymb}
%
% Get bitmap size and resolutions
%\usepackage{bmpsize}
%
% Make tables look nicer
%\usepackage{booktabs}
%
% Compress and sort multi-citations. You _shouldn't_ use this in most modern
% templates since they have their own style, but it can come in handy.
%\usepackage{cite}
%
% Highly customizable enumeration, including inline lists.
%\usepackage[inline]{enumitem}
%
% Enable message sequence charts, which are useful for showing messages sent in
% a particular order (e.g., in network protocols).
%\usepackage{msc}
%
% Enable multiple subfigures within a single figure.
%\usepackage{subfig}
%
% Enable table columns to fill/wrap a specific width. This isn't allowed by ACM
% templates, so your source might get rejected by TAPS if you use it.
%\usepackage{tabularx}
%
% Enable various text symbols (e.g., currency symbols, bullet points, copyright
% symbols, musical notation).
%\usepackage{textcomp}
%
% Make URLs display correctly in citations. The hyphens option allows line
% breaks after hyphens.
%\usepackage[hyphens]{url}
%
% For spacing after user-defined commands and others.
%\usepackage{xspace}
%%%%%%%%%%%%%%%%%%%%%%%%%%%%%%%%%%%%%%%%%%%%%%%%%%%%%%%%%%%%%%%%%%%%%%%%%%%%%%%%


% Acronym definitions go here (requires acronym package). They should be
% organized in order of acronym identifier.

% The identifier used here should have a width that is at least that of the
% widest acronym if generating a list of acronyms used; otherwise, it doesn't
% matter.

% Example acronym definition showing typically used options. Only the first
% line is necessary.
\newacro{as}[AS]{Acronym Society}
\newacroindefinite{as}{an}{an}
\newacroplural{as}[ASes]{Acronym Societies}


% User-defined commands go here.

% GENERAL CUSTOM COMMANDS %%%%%%%%%%%%%%%%%%%%%%%%%%%%%%%%%%%%%%%%%%%%%%%%%%%%%%
% These should work no matter what packages are used (unless the identifiers
% conflict).
%
% Our research tends to quote from interviews a lot, and it can be annoying to
% try and change how all of the quotes are formatted, so set a single format
% here.
\newcommand{\say}[1]{``#1''}
\newcommand{\paraphrase}[1]{[\emph{#1}]}
%
% Similarly, if identifying themes in text, it's helpful to format those in a
% consistent way.
\newcommand{\theme}[1]{\textbf{#1}}
%
% Improve formatting of paragraphs.
%\renewcommand{\paragraph}[1]{\medskip\noindent\textbf{#1}.}
%%%%%%%%%%%%%%%%%%%%%%%%%%%%%%%%%%%%%%%%%%%%%%%%%%%%%%%%%%%%%%%%%%%%%%%%%%%%%%%%

% PACKAGE-SPECIFIC CUSTOMIZATIONS %%%%%%%%%%%%%%%%%%%%%%%%%%%%%%%%%%%%%%%%%%%%%%
% Below are new commands that only work if a specific package is used. For each
% section, begin with a comment line that starts with "Required: ..." to clearly
% mark that it's a package-specific customization.
%
% Required: cite (possibly with noadjust option)
% Fix ugly citation style: [1],[2] --> [1,2].
%\renewcommand{\citepunct}{, }
%\renewcommand{\citedash}{--}
%
% Required: enumitem
% Environment to define commonly-used list types.
% - The default enumerated list has quite a bit of spacing, so make it more
%   compact.
\newenvironment{compactenum}[1][\arabic*.]
{\begin{enumerate}[label={#1}, noitemsep]}
{\end{enumerate}}
% - Same with the itemized list.
\newenvironment{compactitem}%
{\begin{itemize}[noitemsep]}
{\end{itemize}}
% - Customize a list that appears in the flow of a paragraph like (1) one, (2)
%   two, (3) three, (4) four.
\newenvironment{inlineenum}[1][(\arabic*)]
{\begin{enumerate*}[label={#1}, itemjoin={{, }}]}
{\end{enumerate*}}
% - Same as inlineenum, but adds an "and" before the last element to flow more
%   nicely with a sentence, like (1) one, (2) two, and (3) three.
\newenvironment{inparaenum}[1][(\arabic*)]
{\begin{enumerate*}[label={#1}, itemjoin={{, }}, itemjoin*={{, and }}]}
{\end{enumerate*}}
%
% Required: hyperref
% The \autoref command ensures that the proper term is used for the relevant
% reference. Customize the relevant terms for each reference type here.
%\renewcommand*{\sectionautorefname}{Sec.}
%\renewcommand*{\figureautorefname}{Fig.}
%\renewcommand*{\tableautorefname}{Tab.}
% In some templates, subsections are referred to as "Subsection x.y", while
% they may be referred to as "Section x.y" instead. The following line ensures
% that sections and subsections are always referred to with the same identifier.
%\let\subsectionautorefname\sectionautorefname%
%
% Required: msc
% Message sequence charts by default display "msc" in the corner of the chart.
% Turn that off here.
%\setmsckeyword{}
%
% Required: xcolor
% It can be helpful to include formatting for draft text or different authors'
% comments, each with a different text color. Define those here, and also ensure
% that they all get stripped out upon submission.
\ifsubmit%
  \newcommand{\draft}[1]{}
\else
  \newcommand{\draft}[1]{\textcolor{blue}{[Draft text: #1]}}
\fi
%
% Required: xspace
% For some papers, we define new notation that only works in math mode. Ensure
% that these can be easily mentioned in text, with the proper spacing, using the
% following wrapper command.
%\newcommand{\emxs}[1]{\ensuremath{#1}\xspace}
%
%%%%%%%%%%%%%%%%%%%%%%%%%%%%%%%%%%%%%%%%%%%%%%%%%%%%%%%%%%%%%%%%%%%%%%%%%%%%%%%%


% Any miscellaneous preamble commands can be done here.

% The \DocumentMetadata command in the ACM template enables PDF tagging, which
% due to a bug in the acmart class leads to a blank page being created at the
% start of the PDF. Patch the command here to prevent this. This might break PDF
% tagging in places for some types of ACM templates; hopefully it gets fixed
% upstream soon.
\makeatletter
\IfDocumentMetadataT{%
	\pretocmd{\@printtopmatter}{\tagpdfparaOff}{}{\PackageWarning{paper}{Could not patch \string\@printtopmatter}}%
	\apptocmd{\@printtopmatter}{\tagpdfparaOn}{}{\PackageWarning{paper}{Could not patch \string\@printtopmatter}}%
}
\makeatother





\begin{document}

\title{Title Here}

\ifcameraready{}
  % This is an author block - add one for each author.
  \author{Real Name 1}
  \email{Real Email 1}
  \orcid{Real ORCID 1}
  % Below is the affiliation. Multiple authors can share one affiliation;
  % just copy the above lines as many times as needed for authors sharing an
  % affiliation.
  \affiliation{%
    \institution{Real Institution 1}
    \city{Real City 1}
    % State is optional; include it according to the customs of the country in
    % which your institution is located.
    \state{Real State 1}
    \country{Real Country 1}
  }
  % Copied from ACM template.
  % By default, the full list of authors will be used in the page headers.
  % Often, this list is too long, and will overlap other information printed in
  % the page headers. This command allows the author to define a more concise
  % list of authors' names for this purpose.
  %\renewcommand{\shortauthors}{Real Name et al.}
\else
  % Some venues use anonymized author blocks like this rather than the anonymous
  % option in the acmart document class, to ensure that the number of pages is
  % accurate. In that case, make author blocks in the same structure as for real
  % names above.
  \author{Placeholder Name 1}
  \email{Placeholder Email 1}
  \orcid{Placeholder ORCID 1}
  \affiliation{%
    \institution{Placeholder Institution 1}
    \city{Placeholder City 1}
    \state{Placeholder State 1}
    \country{Placeholder Country 1}
  }
  %\renewcommand{\shortauthors}{Placeholder Name et al.}
\fi

% For the ACM format, the abstract comes before the maketitle command.
\begin{abstract}
  A typical abstract \emph{roughly} follows the structure described here,
  but some venues may specify a structure or other constraints like word limits.
  Summarize the topic that the paper is in, leading into the main problem. Then
  summarize the main problem that the paper aims to solve. In one or two
  sentences, describe what you did in the paper. If needed, state the research
  questions here. Perhaps most importantly, \emph{summarize} the findings of the
  paper here; don't just write about what will be presented. Finally, put these
  findings in the context of the problem and of potential future work.
\end{abstract}


% ACM requires information on the relevant CCS concepts and keywords. Include
% that here.
%%
%% The code below is generated by the tool at http://dl.acm.org/ccs.cfm.
%% Please copy and paste the code instead of the example below.
%%
\begin{CCSXML}
<ccs2012>
 <concept>
  <concept_id>00000000.0000000.0000000</concept_id>
  <concept_desc>Do Not Use This Code, Generate the Correct Terms for Your Paper</concept_desc>
  <concept_significance>500</concept_significance>
 </concept>
 <concept>
  <concept_id>00000000.00000000.00000000</concept_id>
  <concept_desc>Do Not Use This Code, Generate the Correct Terms for Your Paper</concept_desc>
  <concept_significance>300</concept_significance>
 </concept>
 <concept>
  <concept_id>00000000.00000000.00000000</concept_id>
  <concept_desc>Do Not Use This Code, Generate the Correct Terms for Your Paper</concept_desc>
  <concept_significance>100</concept_significance>
 </concept>
 <concept>
  <concept_id>00000000.00000000.00000000</concept_id>
  <concept_desc>Do Not Use This Code, Generate the Correct Terms for Your Paper</concept_desc>
  <concept_significance>100</concept_significance>
 </concept>
</ccs2012>
\end{CCSXML}

\ccsdesc[500]{Do Not Use This Code~Generate the Correct Terms for Your Paper}
\ccsdesc[300]{Do Not Use This Code~Generate the Correct Terms for Your Paper}
\ccsdesc{Do Not Use This Code~Generate the Correct Terms for Your Paper}
\ccsdesc[100]{Do Not Use This Code~Generate the Correct Terms for Your Paper}

%%
%% Keywords. The author(s) should pick words that accurately describe
%% the work being presented. Separate the keywords with commas.
\keywords{Do, Not, Us, This, Code, Put, the, Correct, Terms, for,
  Your, Paper}



% ACM also now includes dates of submission/revision/publication. Include those
% here.
% Adjust these dates if needed.
\received{20 February 2007}
\received[revised]{12 March 2009}
\received[accepted]{5 June 2009}


% This is where the title, authors, etc. are actually placed. The location of
% this command, or the command itself, may vary depending on the template.
% Unlike most other templates, we don't also call \thispagestyle{empty} because
% we've redefined \maketitle in a way that doesn't change the page style.
\maketitle

% Place all needed sections, including the abstract, here. In some cases, other
% commands may need to be placed in between sections (e.g., some styles place
% keywords between the abstract and main paper body). There's almost always an
% abstract, intro, and conclusion, so place those here.

\section{Introduction}\label{sec:intro}

% TODO: Expand this section with more detail to help provide structure to those
% newer to writing papers.

Motivate and explain the problem that the paper is trying to solve. If
necessary, explain concepts/terminology and cite sources to back up the claims
you are making. If the paper is tackling a specific part of the problem, be sure
to motivate that part as well.

In order to address this problem, describe what we actually do in the paper.
Summarize the methods we used and explain how the results are useful. Explain
the contributions this paper makes towards addressing the problem. If needed,
clearly state the research questions investigated.


\section{Theoretical Frameworks}\label{sec:theory}

Our research typically relies on established theoretical frameworks as a lens
for collecting and analyzing data. Explain relevant frameworks here.


\section{Methods}\label{sec:methods}

Summarize overall methods here. Typically, it's helpful to describe an overall
approach to the methods first, which might include a discussion of one or more
methodologies and why they're suitable, explicitly connecting the methods to
research questions, etc.

\subsection{Data Collection}\label{sec:methods:collection}

For the kind of work we do, explain the data source, including what data was
collected (e.g., interviews, observations, materials) and from whom. Explain how
subjects were selected. If qualitative data was collected (e.g., through
interviews), explain the experience and perspective of the person who
designed/asked the questions. Be sure identify potential sources of bias (e.g.,
terminology used in questions or the selection of the specific materials).

Ideally, we should explain the specifics of data collection, which might include
the modality of data collection (in person versus remote), questions asked, etc.
If necessary, also describe how participants were selected and why, and evidence
that a sufficiently representative set of participants was selected, and/or that
saturation was reached.

\subsection{Data Analysis}
\label{sec:methods:analysis}

Explain how we analyzed the data. Again, if necessary, include citations or
positionality statements.


\section{Results}\label{sec:results}

Explain results here.


\section{Discussion}\label{sec:discussion}

Describe potential limitations, alternative explanations of findings, what can
be learned, etc.

Sometimes conclusions can be here, or sometimes there may be an additional
section.


%%
%% The acknowledgments section is defined using the "acks" environment
%% (and NOT an unnumbered section). This ensures the proper
%% identification of the section in the article metadata, and the
%% consistent spelling of the heading.
\begin{acks}
  This material is based upon work supported by the XYZ Foundation
  under Award Nos.\@ 123, 456, and 789. The authors would like to thank Alice,
  Bob, and Charlie for help in programming, proofreading, etc.
\end{acks}


% Specify the bibliography style and file here. The style will depend on which
% template you use, but the file (bib.bib in this case) can be left as is,
% unless you rename that file.
\bibliographystyle{ACM-Reference-Format}
% Some venues require you to balance the final two columns, which typically
% consists of references. Include the line below if needed, or comment it out if
% not.
\balance{}
\bibliography{bib}

% Some templates put appendices after the bibliography. If so, feel free to add
% them here. Make sure that the relevant files (e.g., app/foo.tex) exist.
%\appendix
%\input{app/foo}

\end{document}

% Anything here or later will be ignored in the final document, so feel free to
% put notes/etc here.

